% Figura redibujada a partir de "Imagenes/Espectro.PNG" (Tema 1).
% Coordenadas en píxeles del PNG original (698 x 629 px), escaladas para que
% ocupe lo mismo que el PNG en el documento: 210.2 pt x 189.4 pt.
% Escala de energías: x = 33 + 252.9 (E_g - 1.2), con lambda = 1.24/E_g.
\documentclass[12pt,border=0pt]{standalone}
\usepackage[spanish,es-tabla]{babel}
\usepackage[utf8]{inputenc}
\usepackage{amsmath}
\usepackage{xcolor}
\usepackage{tikz}
\usepackage{pgfplots}
\usepackage{circuitikz}
\usetikzlibrary{babel, circuits.ee}
\usetikzlibrary{arrows.meta, decorations.pathreplacing}

\definecolor{applegreen}{rgb}{0.55, 0.71, 0.0}
\definecolor{airforceblue}{rgb}{0.36, 0.54, 0.66}
\definecolor{amber}{rgb}{1.0, 0.49, 0.0}

\newcommand{\diff}[3][1]{%
	\IfBlankTF{#1}{\frac{\mathrm{d}{#2}}{\mathrm{d}{#3}}}%
	{\frac{\mathrm{d}^{#1} {#2}}{\mathrm{d}{#3}^{#1}}}%
}%
\newcommand{\pdiff}[3][1]{%
	\IfBlankTF{#1}{\frac{\partial {#2}}{\partial {#3}}}%
	{\frac{\partial^{#1} {#2}}{\partial {#3}^{#1}}}%
}%

\begin{document}
\begin{tikzpicture}[x=0.30115pt, y=-0.30115pt, font=\scriptsize,
    declare function={xE(\E) = 33 + 252.9*(\E - 1.2); xL(\l) = xE(1.24/\l);}]
    \useasboundingbox (0,0) rectangle (698,629);
    % ---- Respuesta del ojo ----
    \draw[-latex] (153,230) -- (153,85);
    \node[align=center, inner sep=1pt] at (153,48) {Respuesta relativa\\del ojo};
    \draw[blue, ultra thick, domain=153:420, samples=80]
        plot (\x, {230 - 130*exp(-((\x - 290)/36)^2/2)});
    \draw[latex-latex] (249.5,165) -- (330.5,165);
    \draw (290,165) to[out=80, in=180] (345,133);
    \node[anchor=west] at (345,133) {Anchura a media altura};
    \draw[-latex] (420,210) -- (290,210) -- (290,228);
    \node[anchor=west] at (430,210) {$\lambda_\mathrm{m} = 0.555\ \mu$m};
    % ---- Regiones del espectro ----
    \draw (0,230) -- (698,230);
    \draw (0,310) -- (698,310);
    \foreach \x in {225,273,383} \draw (\x,230) -- (\x,270);
    \begin{scope}[font=\fontsize{7}{8}\selectfont]
        \node[anchor=west, inner sep=1pt] at (0,252) {Infrarrojo};
        \node at (189,252) {Rojo};
        \node at (328,252) {Verde};
        \node at (448,252) {Violeta};
        \node at (603,252) {Ultravioleta};
        \node at (183,288) {Naranja};
        \node at (302,288) {Amarillo};
        \node at (410,288) {Azul};
    \end{scope}
    % ---- Gaps de los semiconductores (dx: desplazamiento de la etiqueta) ----
    \foreach \E/\nombre/\dx in {1.12/Si/0, 1.42/GaAs/0, 1.74/CdSe/0, 2.26/GaP/-9,
                                2.42/CdS/9, 3.0/SiC/0, 3.4/GaN/-6, 3.68/ZnS/6} {
        \draw ({xE(\E)},310) -- ({xE(\E)},345);
        \node at ({xE(\E)+\dx},368) {\nombre};
    }
    \draw[decorate, decoration={brace, amplitude=4pt, mirror}]
        ({xE(1.42)-12},392) -- ({xE(2.26)+12},392);
    \node at ({(xE(1.42)+xE(2.26))/2},422) {GaAs$_{1-y}$P$_y$};
    % ---- Escalas de lambda y E_g ----
    \draw ({xE(1.2)},532) -- ({xE(3.6)},532);
    \foreach \E in {1.2,1.4,...,3.61} {
        \draw ({xE(\E)},532) -- ({xE(\E)},545);
        \node at ({xE(\E)},570) {\pgfmathprintnumber[fixed, fixed zerofill, precision=1]{\E}};
    }
    \foreach \E in {1.3,1.5,...,3.51} \draw ({xE(\E)},532) -- ({xE(\E)},539);
    \foreach \l/\dx in {1.0/-7,0.9/5,0.8/0,0.7/0,0.6/0,0.5/0,0.45/0,0.4/0,0.35/0} {
        \draw ({xL(\l)},532) -- ({xL(\l)},519);
        \node[font=\fontsize{7}{8}\selectfont] at ({xL(\l)+\dx},507) {$\l$};
    }
    \foreach \l in {0.95,0.85,0.75,0.65,0.55,0.375} \draw ({xL(\l)},532) -- ({xL(\l)},525);
    \node at (350,472) {$\lambda$ ($\mu$m)};
    \node at (335,603) {$E_\mathrm{g}$ (eV)};
\end{tikzpicture}
\end{document}
